% Part: normal-modal-logic
% Chapter: axioms-systems
% Section: soundness

\documentclass[../../../include/open-logic-section]{subfiles}

\begin{document}

\olfileid{nml}{prf}{snd}

\olsection{سموالے}

!!^a{derivation} نظام هله سم بلل کېږي چې هر هغه څه چې ترې
!!{derive} کېدلے شي، معتبر وي۔ د مودالي نظامونو په پام کښې نيولو سره،
يعنې داسې !!{derivation}s چې د \Ax{K} ترڅنګ د ځينو !!{formula}s
$!A_1$، \dots،~$!A_n$ بېلګې هم پکې کارولے شو، غواړو هر !!{derivable}
فارمول په هر هغه مدل کښې رښتيا وي چې $!A_1$، \dots، $!A_n$ پکې
رښتيا وي۔

\begin{thm}[د سموالي قضيه]\ollabel{thm:soundness}
  که د $!A_1$، \dots، $!A_n$ هره بېلګه په ترتيب سره د مدلونو په
  ټولګيو $\mClass{C}_1$، \dots، $\mClass{C}_n$ کښې معتبره وي،
  نو له $\Log{K}!A_1\dots !A_n
  \Proves !B$ څخه لازمېږي چې $!B$ د مدلونو په ټولګي
  $\mClass{C}_1 \cap \dots \cap \mClass{C}_n$ کښې معتبر وي۔
\end{thm}

\begin{proof}
  د ثبوتونو پر اوږدوالي استقرا کوو۔ د لنډون دپاره $\mClass{C} =
  \mClass{C}_1 \cap \dots \cap \mClass{C}_n$ وټاکئ۔
  \begin{enumerate}
  \item د استقرا بنسټ: که د $!B$ د ثبوت اوږدوالے~$1$ وي، نو دا يا
    د همېش‌رښتيا فارمول يوه بېلګه وي، يا د~\Ax{K} يوه بېلګه،
    \iftag{prvDiamond}{ يا د~\Dual{} يوه بېلګه،}{} يا د
    $!A_1$، \dots،~$!A_n$ له ډلې د يوه فارمول بېلګه وي۔ په لومړي
    حالت کښې $!B$ په $\mClass{C}$ کښې معتبر وي، ځکه د همېش‌رښتيا
    فارمول بېلګې د مدلونو په \emph{هره} ټولګي کښې معتبرې دي؛
    \olref[syn][tau]{prop:valid-taut} وګورئ۔ په دويم حالت کښې هم
    د \olref[syn][sch]{prop:Kvalid}\iftag{prvDiamond}{ او
      \olref[syn][sch]{prop:Dual-valid}}{} له مخې همدا پايله راځي۔
    په درېيم حالت کښې، څرنګه چې $!B$ په $\mClass{C}_i$ کښې معتبر دے
    او $\mClass{C} \subseteq
    \mClass{C}_i$، نو د \olref[syn][val]{prop:subset-class} له مخې
    $!B$ په $\mClass{C}$ کښې هم معتبر دے۔
  \item د استقرا ګام: فرض کړئ $!B$ د $k>1$ اوږدوالي ثبوت لري۔ که
    $!B$ د همېش‌رښتيا فارمول يوه بېلګه يا د $!A_1$،
    \dots، $!A_n$ له ډلې د يوه بېلګه وي، د تېر ګام په شان استدلال
    کوو۔ نو فرض کړئ $!B$ د پخوانيو !!{formula}s $!C \lif !B$ او
    $!C$ څخه د \MP{} په واسطه لاس ته راغلے وي۔ د $!C \lif !B$ او
    $!C$ د ثبوتونو اوږدوالے $<k$ دے، او د استقرا د فرضيې له مخې
    دواړه په~$\mClass{C}$ کښې معتبر دي۔ د
    \olref[syn][sch]{prop:soundMP} له مخې $!B$ هم په $\mClass{C}$
    کښې معتبر دے۔ په پاے کښې فرض کړئ $!B$ له $!C$ څخه د \Nec{} په
    واسطه لاس ته راغلے وي (نو $!B = \Box!C$)۔ د استقرا د فرضيې له
    مخې $!C$ په $\mClass{C}$ کښې معتبر دے، او د
    \olref[syn][val]{prop:Nec-rule} له مخې~$!B$ هم معتبر دے۔
  \end{enumerate}
\end{proof}

\end{document}
